
\clearpage\section{ÁREAS DE ATUAÇÃO}
O \NOMEEGRESSO é o profissional de nível superior com competências e
habilidades para planejar, implementar, administrar, gerenciar, promover e aprimorar com técnica e tecnologia a
automação industrial, assumindo ação empreendedora com consciência de seu papel político, econômico,
social e ambiental. 

É um profissional com formação generalista que atua no controle e automação de
equipamentos, processos, unidades e sistemas de produção, atuando principalmente na interface entre o sistema
produtivo e o sistema gerencial das empresas, planejando, projetando e executando sistemas de controle de processos e
de produção industrial, voltado de modo geral para a automação dos métodos e dos equipamentos industriais. 

Em sua atuação também estuda, projeta e especifica materiais, componentes, dispositivos ou equipamentos
elétricos, eletromecânicos, eletrônicos, magnéticos, ópticos, de instrumentação, de
aquisição de dados e de máquinas elétricas; planeja, projeta, instala, opera e mantém sistemas de
medição e instrumentação eletro-eletrônica, de acionamentos de máquinas, de controle e
automação de processos, de equipamentos dedicados, de comando numérico e de máquinas de
operação autônoma; projeta, instala e mantém robôs, sistemas de manufatura e redes industriais;
coordena e supervisiona equipes de trabalho, realiza estudos de viabilidade técnico-econômica, executa e
fiscaliza obras e serviços técnicos e efetua vistorias, perícias e avaliações, emitindo laudos
e pareceres técnicos. Em todas as suas atividades, considera aspectos referentes à ética, à
segurança, à legislação e aos impactos ambientais, além da preocupação com o uso eficiente das
energias durante o pleno funcionamento de equipamentos e processos fabris.

O aluno egresso do \NOMECURSO tem potencialidade para atuar tanto nas empresas de
engenharia e nas indústrias de produção de equipamentos e software para automação industrial, como nos setores usuários
da automação, podendo sua intervenção acontecer nos seguintes níveis:

\begin{enumerate}
	\item Automatização de processos e sistemas no setor industrial, comercial, residencial e de serviços;
	\item Modernização, otimização do funcionamento e manutenção de unidades de produção automatizadas;
	\item Projeto e integração de sistemas de automação industrial em empresas de engenharia;
	\item Concepção e instalação de unidades de produção automatizadas;
	\item Concepção e fabricação em unidades de produção automatizada;
	\item Desenvolvimento de produtos de instrumentação, controle, operação e supervisão de processos industriais.
	\item Treinamento de recursos humanos em indústrias e instituições de ensino;
	\item Pesquisa científica e tecnológica buscando a criação e desenvolvimento de novas tecnologias.
\end{enumerate}

Neste escopo, fica claro que o Engenheiro de Controle e Automação está habilitado para trabalhar em
concessionárias de energia, automatizando os setores de geração, transmissão ou
distribuição de energia; na automação de indústrias e na automação predial; com
simulação, análise e emulação de grandes sistemas por computador; na fabricação e
aplicação de máquinas e equipamentos elétricos robotizados ou automatizados. Portanto, destacam-se
como possíveis locais de absorção desta mão de obra qualificada, Empresas de Engenharia; Empresas de beneficiamento e
de bebidas; Empresas de linha de montagem industrial; Empresa de geração e distribuição de energia
elétrica; Empresa de prospecção e beneficiamento de petróleo e gás; Empresas de siderurgia,
laminação; Empresas do ramo têxtil e calçadista.

